\scanpage{34}

\uline{Remarque 1} La \uline{propriété $LG$ faible} \uline{équivaut}
à $(LG)_0$.

En effet, il suffit de prouver qu'elle implique que
\[
H^1(X,F)\longrightarrow H^1(X_{/Y},F_{/Y})
\]
est injectif, donc que tt. ext. de $\underline O_X$ par $F$
qui splitte formellement, splitte \unsure{régulièrement|réellement}, ce qui est immédiat.
[Question : si $\alpha$ est un isom pour $i\leq n$,
$\forall$ $F$ loc libre, est-il vrai que $\alpha$ est un mono
pour \deleted{tout} $i=n+1$, $\forall$ $F$ loc libre ?]
\inserted{Oui, prendre une \unsure{résolution} injective,
$0\longrightarrow F\longrightarrow I\longrightarrow G\longrightarrow0$
et \unsure{vérifier que le foncteur} $/Y$, \unsure{et la}
\unsure{limite} des $\unsure{G}=0$]}

\uline{Remarque 2}
Soit $(L_\alpha)$ une famille de Modules loc. libres sur $X$,
telle que $\forall$ module loc. libre sur $X$, est
$\simeq$ quotient d'une somme directe de Modules $\check L_\alpha$,
i.e. $\simeq$ sous-module \inserted{loc. facteur direct}
d'une somme directe de Modules $L_\alpha$.
[Ex : 1) $X$ quasi \deleted{\unsure{projectif}}\inserted{affine},
$(L_\alpha)$ réduit à $\underline O_X$
2) $X$ quasi-proj,
$(L_\alpha)$ formé des $\underline O(m)$, $m\geq m_0$].
Alors $(LG)_n$ équivaut à la condition
\[
H^i(X,L_\alpha)\simeq H^i(X_{/Y},L_{\alpha/Y})
\quad\left|\begin{array}{l}
\text{isom pour }i\leq n\\
\text{inj pour }i=n+1
\end{array}\right.
\]
\deleted{[N.B. \inserted{\unsure{Confiance} :}
\unsure{faut le prouver} \unsure{en mettant une}
condition sur $i=n+1$ !]}.
Dém par récurrence sur $n$, trivial si $n<-1$,
donc $\unsure{n+1<0|n\geq0}$, supposons \deleted{\unsure{on}}
que c'est prouvé

\margin{N.B. Le même argument marche si on
\unsure{ne l'exige que la} condition
$H^{\unsure{n+1|n}}(\unsure{L_\alpha})\longrightarrow
H^{\unsure{n+1|n}}(\unsure{\hat L_\alpha})$
\uline{injectif}.
\deleted{\unsure{Cet|Cat} \unsure{résultat|réserve}}}

\margin{33}
